One concern with the previous algorithms is that they may find a  proportional clustering with poor global objective (e.g., $k$-median), even when exact proportional clusterings with good global objectives exist. For example, suppose $k=2$ and there are two easily defined clusters, containing $40\%$ and $60\%$ of the data respectively. It is possible that Algorithm~\ref{alg:ball} will only open centers inside of the larger cluster. This is proportional, but undesirable from an optimization perspective (note that the ``correct'' clustering of such an example is still proportional). Here, we show how to address this concern by optimizing the $k$-median objective subject to proportionality as a constraint. Later, in Section~\ref{sec:experiments}, we empirically study the tradeoff between the $k$-means objective and proportionality on real data.

We consider the $k$-median and $k$-means objectives to be reasonable measures of the global quality of a solution. We see minimizing the $k$-center objective more as a competing notion of fairness, and so we focus on optimizing the $k$-median objective subject to proportionality.\footnote{A constant approximation algorithm for minimizing the $k$-median objective immediately implies a constant approximation algorithm for minimizing the $k$-means objective by running the algorithm on the squared distances~\cite{mettuPlaxton}.} %We therefore focus on minimizing the $k$-median objective, subject to proportionality as a constraint. 
Minimizing the $k$-median objective without proportionality is a well studied problem in approximation algorithms, and several constant approximations are known~\cite{kMedianLP, kmedians, mettuPlaxton}. %It is NP-hard to approximate to within less than a $1+2/e$ factor~\cite{kMedianApprox}, and the . 
Most of this work is in the model where $\N \subseteq \M$, and we follow suit in this section. %We extend the linear programming approach of~\cite{kMedianLP} 
We show the following.

\begin{theorem}
\label{thm:lp}
Suppose there is a $\rho$-proportional clustering with $k$-median objective $c$. In polynomial time in $m$ and $n$, we can compute a $O(\rho)$-proportional clustering with $k$-median objective at most $8 c$.
\end{theorem}

In particular, we can compute a constant approximate proportional clustering with $k$-median objective at most eight times the minimum $k$-median objective proportional clustering. Note that the exact running time will depend on the algorithm used to solve the linear program. In the remainder of this section, we will sketch the proof of Theorem~\ref{thm:lp}. We begin with the standard linear programming relaxation of the $k$-median minimization problem, and then add a constraint to encode proportionality. The final linear program is shown in Figure~\ref{fig:lp}. Recall that $B(x, \delta) = \{i \in \N: \; d(i, x) \leq \delta\}$.

\begin{figure}
\begin{align}
	\label{obj} & \mbox{Minimize} & \sum_{i \in \N} \sum_{j \in \M} d(i, j) z_{ij} && \\
	\label{con1} & \mbox{Subject to} & \sum_{j \in \M} z_{ij} = 1 &  & \forall i \in \N \\
	\label{con2} && z_{ij} \leq y_j  &  & \forall j \in \M, \forall i \in \N \\
	\label{con3} && \sum_{j \in \M} y_j \leq k && \\
	\label{con4} && \sum_{j' \in B(j, \gamma R_j)} y_{j'} \geq 1 & & \forall j \in \M \\
	\label{con5} && z_{ij}, y_j \in [0, 1] & & \forall j \in \M, \forall i \in \N
\end{align} 
\caption{Proportional $k$-median Linear Program}
\label{fig:lp}
\end{figure}

In the LP, $z_{ij}$ is an indicator variable equal to 1 if $i \in \N$ is matched to $j \in \M$. $y_j$ is an indicator variable equal to $1$ if $j \in X$, i.e., if we want to use center $j \in \M$ in our clustering. Objective~\ref{obj} is the $k$-median objective. Constraint~\ref{con1} requires that every point be matched, and constraint~\ref{con2} only allows a point to be matched to an open center. Constraint~\ref{con3} allows at most $k$ centers to be opened, and constraint~\ref{con5} relaxes the indicator variables to real values between 0 and 1. 

Constraint~\ref{con4} is the new constraint that we introduce. Our crucial lemma argues that constraint~\ref{con4} approximately encodes proportionality. Let $R_j$ be the minimum value such that $|B(j, R_j)| \geq \lceil \frac{n}{k} \rceil$. In other words, $R_j$ is the distance of the $\lceil \frac{n}{k} \rceil$ farthest point in $\N$ from $j$.

\begin{lemma}
\label{lemma:fairnessConstraint}
	Let $X$ be a clustering, and let $\gamma \geq 1$. If \,$\forall j \in \M$ there exists some $x \in X$ such that $d(j, x) \leq \gamma R_j$, then $X$ is $(1+\gamma)$-proportional. If $X$ is $\gamma$-proportional, then $\forall j \in \M$ there exists some $x \in X$ such that $d(j, x) \leq (1+\gamma) R_j$.
\end{lemma}
\begin{proof}
	Suppose that $\forall j \in \M$ there exists some $x \in X$ such that $d(j, x) \leq \gamma R_j$. Suppose for a contradiction that $X$ is not $(1+\gamma)$-proportional. Then there exists $S \subseteq \N$ with $|S| \geq \lceil \frac{n}{k} \rceil$ and $j \in \M$ such that $\forall i \in S, \; (1+\gamma) \cdot d(i, j) < D_i(X)$. By assumption, $\exists x \in X$ such that $d(j, x) \leq \gamma R_j$, so by the triangle inequality $D_i(X) \leq d(i, j) + d(j, x) \leq d(i, j) + \gamma R_j$. Therefore, $\forall i \in S, \; \gamma \cdot d(i, j) < D_i(X) - d(i, j) \leq \gamma R_j$. However, by definition of $R_j$, since $|S| = \lceil \frac{n}{k} \rceil$, there must exist some $i \in S$ such that $d(i, j) \geq R_j$.
	
	Suppose that $X$ is $\gamma$-proportional. Let $j \in \M$. Consider the set $S$ of the closest $\lceil \frac{n}{k} \rceil$ points in $\N$ to $j$. By definition of proportionality $\exists i \in S$ and $x \in X$ such that $\gamma d(i, j) \geq d(i, x)$. Therefore, by the triangle inequality, $d(j, x) \leq d(i, j) + d(i, x) \leq (1+\gamma) d(i, j)$. By definition of $S$, $d(i, j) \leq R_j$, so there exists $x \in X$ such that $d(j, x) \leq (1+\gamma) R_j$.  
\end{proof}

Now, suppose there is a $\rho$-proportional clustering $X$ with $k$-median objective $c$. Then we write the linear program shown in Figure~\ref{fig:lp} with $\gamma = \rho+1$ in constraint~\ref{con4}. Lemma~\ref{lemma:fairnessConstraint} guarantees that $X$ is feasible for the resulting linear program, so the fractional solution has $k$-median objective at most $c$.  We then round the resulting fractional solution. In \cite{kMedianLP}, the authors give a rounding algorithm for the the linear program in Figure~\ref{fig:lp} without Constraint~\ref{con4}. %They show that their rounding algorithm is 8-approximate on the $k$-median objective with respect to the optimal fractional solution. 
We show that a slight modification to this rounding algorithm also %to the fractional optimum of the linear program in Figure~\ref{fig:lp} (note that the variables are the same) 
preserves Constraint~\ref{con4} to a constant approximation.

\begin{lemma} (Proved in Section~\ref{sec:lemmaProof})
\label{lemma:rounding}
Let $\{y_j\}, \{z_{ij}\}$ be a fractional solution to the linear program in Figure~\ref{fig:lp}. Then there is an integer solution $\{\hat{y_j}\}, \{ \hat{z_{ij}}\}$ that is an $8$-approximation to the objective, and that opens $k$ centers. Furthermore, for all $j \in \M$, $\sum_{j' \in B(j, \, 27 \gamma R_j)} \hat{y_{j'}} \geq 1$.
\end{lemma}

%The proof of Lemma~\ref{lemma:rounding} involves the technical details of the rounding from~\cite{kMedianLP}, and is presented with the supplementary materials. 
The proof of Lemma~\ref{lemma:rounding} is in parts and involves the technical details of~\cite{kMedianLP}. We provide the proof in section~\ref{sec:lemmaProof}, but first we complete the overall proof of Theorem~\ref{thm:lp}. Given Lemma~\ref{lemma:rounding}, applying Lemma~\ref{lemma:fairnessConstraint} again implies that the result of the rounding is $(27(1+\rho)+1)$-proportional, since we set $\gamma = 1+\rho$. Since the $k$-median objective of the fractional solution is at most $c$, the fact that the $k$-median objective of the rounded solution is at most $8c$ follows directly from the proof from~\cite{kMedianLP}. We note that the constant factor of 27 can be improved to 13 in the special case where $\N = \M$. Interestingly, the ostensibly similar primal-dual approach of~\cite{kMedianPrimalDual} does not appear amenable to the added constraint of proportionality (in particular, the reduction to facility location from~\cite{kMedianPrimalDual} is no longer straightforward).    

%%%%%%%%%%%%%%%%%%%%%
%%%%%%%%%%%%%%%%%%%%%
%%%%%%%%%%%%%%%%%%%%%

\subsection{Proof of Lemma~\ref{lemma:rounding}}
\label{sec:lemmaProof}

First, we give a brief overview of the method from \cite{kMedianLP}. Note that our goal in this argument is to show that the new constraint we added (constraint~\ref{con4}) is approximately satisfied after this rounding. The authors work in a setting where points can have \textit{demand}. In our setting, this just corresponds to points in $\N$ having a demand of 1, and moving or consolidating demand can be thought of as changing the instance by moving points in $\N$. Note that the original linear program requires $\N \subseteq \M$, and we follow suit in this proof. 

Given a fractional solution to the linear program in Figure~\ref{fig:lp} as $\{y_j\}, \{z_{ij}\}$, let $\bar{C_i} = \sum_{j \in \M} d(i, j) z_{ij}$. That is, $\bar{C_i}$ is the contribution of point $i$ to the $k$-median objective in the fractional optimum.
\begin{itemize}
  \item Step 1: Consolidate all demands $t_i$ to obtain $\{t'_i\}$ such that for all $i, j \in \M$ with $t'_i > 0, t'_j > 0$, they must be sufficiently far away such that $c_{ij}>4\max(\bar{C_i},\bar{C_j})$. Let $\M'$ be the set of centers with positive demand after this step, i.e. $\M' = \{j \in \M: t'_j > 0\}$.
  \item Step 2a: Consolidate open centers by moving each center not in $\M'$ to the nearest center in $\M'$. This gives a new solution $\{y'_j\}, \{z'_{ij}\}$ with $y'_j \geq \frac{1}{2}$ if $j \in \M'$ and $y'_j = 0$ if $j \notin \M'$. We call this a $\frac{1}{2}$-restricted solution.
  \item Step 2b: Modify the solution further to obtain $\{ \bar{y_j} \}, \{ \bar{z_{ij}} \}$ with $\bar{y_j} \in \{ \frac{1}{2}, 1 \}$ if $j \in \M'$ and $\bar{y_j} = 0$ if $j \notin \M'$. We call this a $\lbrace \frac{1}{2},1 \rbrace$-integral solution.
  \item Step 3: Round $\{ \bar{y_j} \}, \{ \bar{z_{ij}} \}$ to obtain an integer solution $\{ \hat{y_j} \}, \{ \hat{z_{ij}} \}$.
\end{itemize}

We introduce constraint~\ref{con4} in our linear program, and make a small modification to Step 2b as described in sections below.

\subsubsection{Step 1: Consolidating Demands}
%\begin{lemma}
%\label{lemma:demandIn4Ri}
%    Fix $i \in \M$. There exists a location $j\in \M'$ such that $d(i, j)\leq 4 \gamma R_i$.
%\end{lemma}
%\begin{proof}
%    Let $j$ be the location to which the demand for $i$ was moved in Step 1. Note that by definition of Step 1, $c_{ij} \leq 4 \bar{C_i}$. By constraint~\ref{con4}, we know that the demand of $i$ could be completely satisfied by centers inside of $B(i, \gamma R_i)$, so $\bar{C_i} \leq \gamma R_i \implies c_{ij} \leq 4 \gamma R_i$.
%\end{proof}
Our first observation is that during the demand consolidation, it cannot be the case that all of the demand within a given $B(i, \gamma R_i)$ from constraint~\ref{con4} is moved arbitrarily far away.

\begin{lemma}
\label{lemma:demandMovedTo5Ri}
    Fix $i \in \M$. For each $j \in B(i, \gamma R_i)$ that had its demand moved to $j' \in \M'$, $d(i, j') \leq 9 \gamma R_i$.
\end{lemma}
\begin{proof}
    %By Lemma~\ref{lemma:demandIn4Ri}, $\exists j''\in \M'$ with $d(i, j'') \leq 4 \gamma R_i$. Since $j'$ is the location in $\M'$ closest to $j$, $d(j, j') \leq d(j, j'') \leq d(i, j) + d(i, j'') \leq \gamma R_i + 4 \gamma R_i = 5 \gamma R_i$.
    Let $j \in B(i, \gamma R_i)$. Let $j'$ be the location to which the demand for $j$ was moved in Step 1. Note that Step 1 is designed so that if demand at $j$ is moved to another point $j'$, then $\bar{C_{j'}} \leq \bar{C_j}$, and $c_{jj'} \leq 4\max(\bar{C_j},\bar{C_{j'}}) = 4 \bar{C_j}$. By constraint~\ref{con4}, we know that the demand of $j$ could be completely satisfied by centers fractionally opened inside of $B(i, \gamma R_i)$, so $\bar{C_j} \leq 2 \gamma R_i \implies c_{jj'} \leq 8 \gamma R_i$. Since $j \in B(i, \gamma R_i)$, it follows that $d(i, j') \leq 9 R_i$.
\end{proof}

\subsubsection{Step 2: Consolidating Centers}
Next, we argue that the consolidation of fractional centers in Step 2 approximately preserves constraint~\ref{con4}.

\begin{lemma}
\label{lemma:centerIn5Ri}
    After Step 2a, $\forall j \in \M$, $\sum_{j' \in B(j, 9 \gamma R_j)} y'_{j'} \geq 1$.
\end{lemma}
\begin{proof}
    Consider each location $j' \in B(j, \gamma R_j)$. Let $j''$ be the location in $\M'$ closest to $j'$. By Lemma~\ref{lemma:demandMovedTo5Ri}, there exists such a $j''$ with $d(j, j'') \leq 9 \gamma R_j$, so $j'' \in B(j, 9 \gamma R_j)$. Step 2a will move the fractional center at $j'$ to $j''$, so $y'_{j''} \geq \min(1, y_{j'} + y_{j''}) \geq y_{j'}$. Summing over all $j' \in B(j, \gamma R_j)$, we have $\sum_{j'' \in B(j, 9 \gamma R_j)} y'_{j''} \geq \sum_{j' \in B(j, \gamma R_j)} y_j' \geq 1$ by constraint~\ref{con4}.
\end{proof}

For our algorithm, we slightly change Step 2b to the following: Let $\M'' = \lbrace j\in \M': \; y'_j < 1 \rbrace$, $m'=|\M'|$ and $m''=|\M''|$. Sort the locations $j \in \M''$ in decreasing order of $t'_j d(s(j), j)$, where $s(j)$ is the location in $\M'$ closest to $j$ (other than $j$). Set $\bar{y_j} = 1$ for the first $2k - 2m' + m''$ locations in $\M''$ or if $j \in \M'\backslash \M''$, and $\bar{y_j} = \frac{1}{2}$ otherwise. In other words, the only difference from the original Step 2b in \cite{kMedianLP} is that points with integral $y'_j = 1$ will not participate in the sorting; we simply set $\bar{y_j} = 1$ for such points, and then perform the standard rounding on $M''$. \cite{kMedianLP} use the following statement in their proof, and we show that it still holds after our modification to Step 2b.

\begin{lemma}
\label{lemma:roundingCost}
    For any $\frac{1}{2}$-restricted solution $\{y'_j\}, \{z'_{ij}\}$, the modified Step 2b givs a $\lbrace \frac{1}{2},1 \rbrace$-integral solution with no greater cost.
\end{lemma}
\begin{proof}
    From Lemma 7 in \cite{kMedianLP}, the cost of the $\frac{1}{2}$-restricted solution $\{y'_j\}, \{z'_{ij}\}$ is
    \begin{flalign*}
        &&&&&\phantom{{}={}} \sum_{j \in \M'} t'_j \left( 1 - y'_j \right) d(s(j), j)&&\\
        &&&&&=\sum_{j \in \M''} t'_j \left( 1 - y'_j \right) d(s(j), j)&&\\
        &&&&&=\sum_{j \in \M''} t'_j d(s(j), j))-\sum_{j \in \M''} t'_j d(s(j), j) y'_j.&&
    \end{flalign*}
    where the second line follows because $y'_j = 1 \; \forall j\in \M' \backslash \M''$. Our algorithm in Step 2b maximizes $\sum_{j\in \M''} t'_j d(s(j), j) \bar{y_j}$ for the given set of $t'_j d(s(j), j)$, hence achieves a cost at most that of $\{y'_j\}, \{z'_{ij}\}$.
\end{proof}

\begin{lemma}
\label{lemma:twoHalfCenters}
    After Step 2b, $\forall j \in \M$, there is either at least one $j' \in B(j, 9 \gamma R_j)$ with $\bar{y_{j'}} = 1$ or at least two $j' \in B(j, 9 \gamma R_j)$ with $\bar{y_{j'}} \geq \frac{1}{2}$.
\end{lemma}
\begin{proof}
    Given Lemma~\ref{lemma:centerIn5Ri} and the constraints on $y'_i$ after Step 2a, there must be at least one $j' \in B(j, 9 \gamma R_j)$ with positive demand. If there is exactly one such $j'$, Lemma 3 is equivalent to $y'_{j'} = 1$ and Step 2b will ensure $\bar{y_{j'}} = 1$. If there are at least two such $j'$, all of them will have $\bar{y_{j'}} \geq \frac{1}{2}$ after Step 2b.
\end{proof}

\subsubsection{Step 3: Rounding an Integer Solution}
\cite{kMedianLP} gives two rounding schemes, and we use the first one that at most doubles the cost. The important observation is that any center with $\bar{y_j} = 1$ will be opened in the integral solution (that is, if $\bar{y_j} = 1$ then $\hat{y_j}=1$, and for any center with $\bar{y_j}=\frac{1}{2}$, either $j$ itself or another center in $\M'$ closest to $j$ will be opened. This allows us to complete our argument. %This gives the following lemma:
%\begin{lemma}
%\label{lemma:fullCenter9Ri}
 %   Fix $i \in \M$ and let $j \in \M'$ with $d(i, j) \leq 9 \gamma R_i$. (Lemma~\ref{lemma:demandIn4Ri} guarantees the existence of $j$.) Then there exists $j' \in \M'$ in the integral solution $\{ \hat{y_j} \}, \{ \hat{z_{ij}} \}$ with $\hat{ y_{j'} } = 1$ and $d(j, j') \leq 9 \gamma R_i$.
%\end{lemma}
%\begin{proof}
%    After Step 2b, $\bar{y_j} \in \{ \frac{1}{2}, 1 \}$. If $\bar{y_j} = 1$, center $j$ itself must be opened in the integral solution. If $\bar{y_j} = \frac{1}{2}$, by Lemma~\ref{lemma:demandMovedTo5Ri} and Lemma~\ref{lemma:twoHalfCenters}, there must be another center $j' \in B(i, 5 \gamma R_i)$ with $\bar{y_{j'}} \geq \frac{1}{2}$. Let $s(j)$ be the center in $\M'$ closest to $j$, we have $d(j, s(j)) \leq d(j, j') \leq d(i, j) + d(i, j') \leq 4 \gamma R_i + 5 \gamma R_i = 9 \gamma R_i$. Since either $j$ or $s(j)$ must be opened, there exists an open center at most distance $9 \gamma R_i$ away from $j$.
%\end{proof}

\begin{lemma}
    For all $j \in \M$, $\sum_{j' \in B(j, \, 27 \gamma R_j)} \hat{y_{j'}} \geq 1$.
\end{lemma}
\begin{proof}
    By Lemma~\ref{lemma:twoHalfCenters}, there are two cases: either there is some $j' \in B(j, 9 \gamma R_j)$ with $\bar{y_{j'}} = 1$ or there are at least two $j' \in B(j, 9 \gamma R_j)$ with $\bar{y_{j'}} \geq \frac{1}{2}$. In the first case, $\hat{y_{j'}} = \bar{y_{j'}} = 1$, so the lemma statement clearly holds. In the second case, we are guaranteed that for each point $j'$, either we set $\hat{y_{j'}} = 1$ (in which case the Lemma holds) or we set $\hat{y_{j''}} = 1$ where $j''$ is the closest other at least partially open center. But since in this case there are \textit{two} points in $B(j, 9 \gamma R_j)$ partially open, their pairwise distance is at most $18 \gamma R_j$, so $d(j', j'') \leq 18 \gamma R_j \implies d(j, j'') \leq 27 \gamma R_j.$
%By Lemma~\ref{lemma:demandIn4Ri}, there must be a partially opened center $j' \in \M'$ with $\bar{d(j, j')} \leq 4 \gamma R_j$ after Step 2b. Lemma~\ref{lemma:fullCenter9Ri} then gives an open center $j''$ in the integral solution with $d(j', j'') \leq 9 \gamma R_j$. Therefore, $d(j, j'') \leq d(j, j') + d(j', j'') \leq 13 \gamma R_j$, so $\sum_{i \in B(j, \, 13 \gamma R_j)} \hat{y_i} \geq \hat{y_{j''}} = 1.$
\end{proof}

This concludes the proof of Lemma~\ref{lemma:rounding}. The fact that $\{ \hat{y_j} \}, \{ \hat{z_{ij}} \}$ is an 8-approximation of the objective follows immediately from the proof from \cite{kMedianLP} given Lemma~\ref{lemma:roundingCost}, as all other constraints are still satisfied. Finally, we note that the constant factor of 27 can be tightened to 13 in the special case where $\N = \M$. The argument is essentially the same. The crucial improvement comes from the guarantee that $\forall i \in \M$, there is demand at the \textit{center} of each $B(i, \gamma R_i)$. Tracking this demand throughout the rounding leads to the tightened result.

%The formal proof is provided in the supplementary material. Let $\bar{C_i} = \sum_{j \in \M} d(i, j) z_{i, j}$. That is, $\bar{C_i}$ is the contribution of point $i$ to the $k$-median objective in the fractional optimum. The first step of the rounding consolidates any points that are ``close'' to one another with respect to $C_{ij}$. This is done by introducing ``demands'' (i.e., if two points are consolidated to one point, the demand of the new point will be 2). After this step, for all points $i, i'$ with positive demand, $d(i, i') > 4 \max(\bar{C_i}, \bar{C_{i'}})$. We show that after this consolidation step, the demand from points in $B(j, R_j)$ remain within a constant factor of $R_j$ from $j$. In other words, the points that are close to a given center do not move too far away from that center in the demand consolidation step.

%Subsequently, the rounding simplifies the selection of fractionally open centers by consolidating nearby fractional centers and ensuring that they are either integrally selected or selected to exactly $1/2$. This is done in a way that minimizes the $k$-median objective. We show that after this consolidation step, for every $j \in \M$, either there is some $j'$ with $d(j, j') \leq 9R_j$ and $y_{j'} = 1$, or there are at least two centers $j', j'' \in \M$ within $9R_j$ of $j$ such that $y_{j'} \geq 1/2$ and $y_{j''} \geq 1/2$. That is Constraint~\ref{con4} is either integrally satisfied to within a 5-approximation, or there are at least two centers selected to at least $1/2$ that are fractionally satisfying the constraint.

%Finally, the rounding selects every center that is integrally selected, and ensures that for every center that is selected to $1/2$, either it or the closest center selected to at least $1/2$ will be integrally selected. Therefore, it is straightforward to argue that Constraint~\ref{con5} is satisfied in the final integral solution to within an approximation of 27. Note that the guarantee can be strengthened to $13$-proportionality in the important special case where $\N = \M$, that is, when there is only a single set of points given as input. This follows from a more careful analysis that specifically considers the demand point at the center Constraint~\ref{con4}.




%\begin{theorem}{\cite{kMedianLP}}
%\label{thm:Charikar}
%	Given a fractional clustering that minimizes the $k$-median objective, one can round to create an integral clustering that is an 8-approximation to the optimal $k$-median objective. 
%\end{theorem}










